\documentclass[11pt]{article}
\usepackage[T1]{fontenc}
\usepackage{amsmath,amssymb,amsthm}
\usepackage[margin=1in]{geometry}
\usepackage{booktabs}
\usepackage{microtype}
\usepackage[hidelinks]{hyperref}

\newtheorem{theorem}{Theorem}
\newtheorem{lemma}[theorem]{Lemma}
\newcommand{\F}{\mathbb F}
\newcommand{\dminus}{\mathbin{\dot{-}}}

\title{The Exact Prime Threshold for a Short Arithmetic\
Point-Count Formula for Weierstrass Cubics}
\author{Xiaojian Zeng\\
PBC School of Finance, Tsinghua University, 100190\\
\texttt{xiaojianzeng@foxmail.com}}
\date{}

\begin{document}
\maketitle

\begin{abstract}
Dumitru and Prunescu proved that their short arithmetic term gives the
affine point count of every Weierstrass cubic over \(\F_p\) when \(p\geq17\).
We show that this lower bound is exact for the corresponding universal
statement.  At \(p=5,7,11,13\), respectively, the nonsingular pairs
\((A,B)=(1,1),(0,1),(1,3),(0,4)\) give the actual affine counts
\(8,11,17,20\), whereas the arithmetic term gives \(3,4,6,7\).
The evaluations use integer coefficient extraction and complete quadratic
residue counts.  Consequently, for every odd prime \(p>3\), the term is
correct for all \(A,B\in\F_p\) if and only if \(p\geq17\).
\end{abstract}

\section{Statement and context}

For nonnegative integers, write \(u\dminus v=\max\{u-v,0\}\).  Let \(p>3\)
be an odd prime and represent \(A,B\in\F_p\) by their canonical integers in
\(\{0,\ldots,p-1\}\).  Put
\[
 e=\frac{p-1}{2},\qquad q=2^{p^2},
\]
and define
\begin{align}
 d_p(A,B)
   &=\left(\left\lfloor
       \frac{(q^3+Aq+B)^e}{q^{p-1}}
      \right\rfloor\bmod q\right),
 &r_p(A,B)&=d_p(A,B)\bmod p,                                      \label{eq:r}\\
 \widehat N_p(A,B)
   &=(p\dminus r_p)+p\Bigl(1\dminus
      \bigl((4p+1)\dminus r_p^2\bigr)\Bigr).                     \label{eq:nhat}
\end{align}
These are equations (4) and (5), with the notation compressed, of Dumitru
and Prunescu~\cite{DumitruPrunescu2026}.  Let
\[
 N_p(A,B)=\#\{(x,y)\in\F_p^2:y^2=x^3+Ax+B\};
\]
this is an affine count, so no point at infinity is included.

\begin{theorem}\label{thm:main}
For every odd prime \(p>3\),
\[
 \bigl(\widehat N_p(A,B)=N_p(A,B)\text{ for every }A,B\in\F_p\bigr)
 \quad\Longleftrightarrow\quad p\geq17.
\]
\end{theorem}

The sufficient direction---the identity for all parameters at every prime
\(p\geq17\)---is precisely Theorem~3 of
\cite{DumitruPrunescu2026}, including singular cubics.  Its recentering step
uses the Hasse bound, originating in Hasse's work on elliptic function
fields~\cite{Hasse1936}.  The contribution here is the converse: four exact
evaluations exhaust the excluded prime range and show that the published
threshold cannot be lowered for the universal assertion.  The biconditional
formulation and these four witnesses are not claims made in
\cite{DumitruPrunescu2026}.

There is a recent, methodologically different example of completing the
boundary left by a prime-field theorem in
Zeng, Liu, and Ratnavelu~\cite{ZengLiuRatnavelu2026}; there the boundary is
settled by a structural equality case rather than by finite counterexamples.
That work is contextual and is not an input to the proof below.

\section{Exact lower-prime evaluations}

Write
\[
 (x^3+Ax+B)^e=\sum_m c_mx^m.
\]
If a monomial contributing to \(c_{p-1}\) chooses \(i\) copies of \(x^3\),
\(j\) copies of \(Ax\), and \(k\) copies of \(B\), then
\begin{equation}
 i+j+k=e,\qquad 3i+j=p-1,                                      \label{eq:constraints}
\end{equation}
and its value is
\begin{equation}
 \frac{e!}{i!j!k!}A^jB^k.                                     \label{eq:multinomial}
\end{equation}
For the four pairs in the abstract, the nonzero terms are as follows.

\begin{center}
\begin{tabular}{@{}cccccc@{}}
\toprule
\(p\) & \(A\) & \(B\) & nonzero \((i,j,k)\) & \(c_{p-1}\) & \(r_p\) \\
\midrule
5  & 1 & 1 & \((1,1,0)\) & \(2\)   & 2 \\
7  & 0 & 1 & \((2,0,1)\) & \(3\)   & 3 \\
11 & 1 & 3 & \((3,1,1)\) & \(60\)  & 5 \\
13 & 0 & 4 & \((4,0,2)\) & \(240\) & 6 \\
\bottomrule
\end{tabular}
\end{center}

These coefficients are the literal digits \(d_p(A,B)\), not merely values
congruent to them modulo \(p\).  Indeed,
\[
 (q^3+Aq+B)^e=\sum_m c_mq^m.
\]
All coefficients are nonnegative, and their sums in the four rows are
\[
 (1+A+B)^e=9,\ 8,\ 3125,\ 15625,
\]
respectively.  These are smaller than the corresponding radices
\(2^{25},2^{49},2^{121},2^{169}\).  Thus there is no carry between base-
\(q\) digits, and the division and remainder in \eqref{eq:r} extract exactly
\(c_{p-1}=2,3,60,240\).  Reduction modulo \(p\) gives the displayed values of
\(r_p\).

It remains to count the points without approximation.  For
\(f(x)=x^3+Ax+B\), direct evaluation at every \(x\in\F_p\) gives
{\small
\begin{align*}
p=5:  &&(f(0),\ldots,f(4))&=(1,3,1,1,4),\\
p=7:  &&(f(0),\ldots,f(6))&=(1,2,2,0,2,0,0),\\
p=11: &&(f(0),\ldots,f(10))&=(3,5,2,0,5,1,5,1,6,4,1),\\
p=13: &&(f(0),\ldots,f(12))&=(4,5,12,5,3,12,12,9,9,5,3,9,3).
\end{align*}}
The nonzero square sets and the resulting complete fibre counts are

\begin{center}
\begin{tabular}{@{}clccccc@{}}
\toprule
\(p\) & nonzero squares & zeros & square values & nonsquares
      & \(N_p\) & \(\widehat N_p\) \\
\midrule
5  & \(\{1,4\}\)             & 0 & 4  & 1 & 8  & 3 \\
7  & \(\{1,2,4\}\)           & 3 & 4  & 0 & 11 & 4 \\
11 & \(\{1,3,4,5,9\}\)       & 1 & 8  & 2 & 17 & 6 \\
13 & \(\{1,3,4,9,10,12\}\)   & 0 & 10 & 3 & 20 & 7 \\
\bottomrule
\end{tabular}
\end{center}

Here a zero value has one square root, a nonzero square has two, and a
nonsquare has none, so the sixth column is \(\text{zeros}+2\cdot
\text{square values}\).  In all four rows \(r_p^2\leq4p\).  Hence the second
term in \eqref{eq:nhat} vanishes and
\(\widehat N_p=p-r_p=3,4,6,7\), proving each mismatch in the last two
columns.  Finally,
\[
 4A^3+27B^2\pmod p=1,6,5,3,
\]
respectively, so all four cubics are nonsingular.

\begin{proof}[Proof of Theorem~\ref{thm:main}]
For every prime \(p\geq17\), Theorem~3 of
\cite{DumitruPrunescu2026} gives \(\widehat N_p(A,B)=N_p(A,B)\) for every
parameter pair, including pairs defining singular cubics.  For orientation,
their proof shows that \(r_p\) is the residue modulo \(p\) of
\(a_p=p-N_p\) and that \(|a_p|\leq2\sqrt p\); when \(p\geq17\), exactly one
of \(r_p\) and \(r_p-p\) lies in that interval, and the test
\(r_p^2\leq4p\) in \eqref{eq:nhat} selects it.

The primes strictly between \(3\) and \(17\) are exactly
\(5,7,11,13\).  The calculations above give one pair \((A,B)\) at each of
these primes for which \(\widehat N_p(A,B)\ne N_p(A,B)\).  Thus the
universally quantified identity fails at every admissible prime below 17 and
holds at every admissible prime at least 17.
\end{proof}

\section{Reproducibility, scope, and assistance}

The four rows were checked with exact integer arithmetic in two separately
written Python implementations.  One evaluates \eqref{eq:r} literally,
expands the polynomial by convolution, and enumerates all \(p^2\) affine
pairs.  The other uses \eqref{eq:multinomial} and counts square-map fibres.
Both reproduce every digit, residue, point count, value of \eqref{eq:nhat},
and discriminant displayed above.  No floating-point computation or random
search enters the result; the tables also provide a self-contained manual
certificate.

This theorem concerns the universal formula over prime fields only.  It does
not classify the parameter pairs for which the term happens to work below
17, and it makes no assertion about composite moduli or extension fields.
The literature comparison was bounded to the cited sources, so no claim of
priority is made for the short sharpness corollary.

OpenAI Codex assisted with organizing the audited argument, checking the
LaTeX presentation, and locating bibliographic records.  The mathematical
claims were checked against the cited primary sources and the exact
implementations described above; language-model output is not used as proof.

\bibliographystyle{plainurl}
\bibliography{references}

\end{document}
